\section{Coarse assembly and locally finite classes}\label{sec:conventions}

We collect the functorial facts needed to carry a family of boundary classes
through the construction. All spaces below are proper metric spaces, and all
maps used to push forward locally finite K-homology are continuous and proper.
A map between discrete metric spaces is \emph{coarse} if it is proper and
bornologous: for every $R$ there is $S$ such that points at distance at most
$R$ have images at distance at most $S$.

\subsection{Rips complexes and their coarse structure}
For a uniformly discrete bounded-geometry space $Y$, the complex $P_r(Y)$ is
locally finite and has uniformly bounded dimension at each fixed $r$.
We use its usual simplicial topology. A concrete proper metric compatible
with that topology and its coarse structure is convenient when the complex
has disconnected components. Regard a point as a finitely supported
probability measure on $Y$ whose support has diameter at most $r$, and use
the transportation metric
\begin{equation}\label{eq:wasserstein}
 d_r(\nu,\eta)=\inf_{\pi}\sum_{x,y\in Y}d_Y(x,y)\pi(x,y),
\end{equation}
where $\pi$ ranges over the couplings of $\nu$ and $\eta$.
On each finite subcomplex this metric gives the usual topology, equivalently
the topology from spherical simplices. The inclusion $y\mapsto\delta_y$ is
isometric, and every point is within distance $r$ of every vertex in its
support. Moreover, a closed ball about $\delta_y$ of radius $R$ has all its supports
in the finite set $B_Y(y,R+r)$. Indeed, for $x,z\in\supp(\nu)$,
$d_Y(y,x)\le d_Y(y,z)+r$; averaging over $z$ gives
$d_Y(y,x)\le d_r(\delta_y,\nu)+r$. The closed ball is a closed subset
of a finite subcomplex and is compact. Local finiteness gives the asserted
topology globally.

In particular, choosing a support vertex gives a coarse inverse to the
vertex inclusion. Distances between chosen support vertices and between
the corresponding points differ by at most $2r$. Thus
\eqref{eq:wasserstein} realizes the usual coarse structure determined by
supports; it also supplies finite distances between disconnected simplicial
components.

\paragraph{Comparison with spherical metrics.}
One can compare this metric directly with the standard semi-spherical extension:
allow paths in spherical simplices, normalized to have edge length one,
and jumps between vertices of cost $d_Y$. Write $d_{\mathrm{ss}}$ for the
infimum of their lengths. At fixed $r$, simplex dimensions have a common
bound. On each simplex, the transportation distance is bounded by a
constant depending only on $r$ times its spherical distance; this follows
by comparing barycentric coordinates in the bounded dimensions and using
$d_r\le (r/2)\norm{\nu-\eta}_1$. The same inequality, with constant one,
holds for vertex jumps. Conversely, one may travel to a support vertex,
make one vertex jump, and travel to the endpoint. Thus, for a finite
constant $c_r$,
\[
 d_r\le c_r d_{\mathrm{ss}},\qquad d_{\mathrm{ss}}\le d_r+2.
\]
The second inequality uses that the minimum distance between the two
supports is at most $d_r$. The inequalities give the same coarse structure;
the finite-subcomplex argument above gives the usual simplicial topology.
These choices give the standard Rips coarsening and Roe-algebra
identifications \cite[Definition~7.2.8, Proposition~7.2.11 and Theorem~7.2.16]{WY}.
The inclusions for increasing
$r$ are continuous proper coarse maps. A proper Lipschitz map of vertex
spaces induces the corresponding proper simplicial maps of Rips complexes.
The support-vertex selections are used for the Roe-algebra K-theory
identifications; locally finite K-homology is pushed forward by the
continuous proper simplicial maps. The selector compositions before and
after an inclusion differ by a bounded distance, since their values lie
in the same bounded support. This also applies to a proper Lipschitz map
of vertex spaces, with the bound given by the diameter of its image support.
Thus the Roe K-theory identifications respect
these maps and the passage to larger scales.

\subsection{Roe algebras and naturality}
On an ample module over a proper metric space, an operator is locally compact
if multiplication on either side by any function in $C_0$ makes it compact.
Its propagation is at most $R$ if its matrix coefficients between sets at
distance greater than $R$ vanish. For a discrete bounded-geometry space $Y$
on $\ell^2(Y)\otimes\HH_0$, a finite-propagation operator is locally compact
exactly when each of its matrix entries in $\mathcal B(\HH_0)$ is compact.
Indeed, the propagation bound leaves only finitely many nonzero entries in
each row and column; finite-support functions give compact products, and
$C_0(Y)$ follows by uniform approximation.

The locally finite analytic index at each Rips scale takes values in the
Roe algebra of that complex. Coarse equivalence with the vertex space
identifies its K-theory with $K_*(C^*(Y))$. The compatible indices define
$\mu_Y$. We use the usual reduced Roe index and its naturality
\cite[Sections~5.1 and~6.7, Lemma~7.1.3 and Theorem~7.2.16]{WY}, with the corrected localization
comparison in the authors' errata \cite[items~30--31]{WYerrata}.
All spaces used in this application are nonempty, and their ample modules
may be chosen with infinite multiplicity. In particular, for a proper
Lipschitz inclusion $f:Y\to Z$ and its covering isometry
$F(\delta_y\otimes v)=\delta_{f(y)}\otimes v$,
\begin{equation}\label{eq:naturality}
 \mu_Z\circ f_*=(\Ad F)_*\circ\mu_Y.
\end{equation}
On vertex Roe algebras $\Ad F$ is an actual homomorphism; the induced map on
K-theory is independent of the covering choices. All the maps to which we
apply \eqref{eq:naturality} will be inclusions of this form.

\subsection{Families and winding pairings}
For a countable topological disjoint union of second-countable locally compact
spaces $Y_j$, all the algebras $C_0(Y_j)$ and their countable direct sum are
separable. Countable additivity in the first variable of $\KK$
\cite[Theorem~19.7.1 and Section~19.8.10]{Blackadar} gives
\begin{equation}\label{eq:additivity}
 K_i^{\lf}\left(\coprod_jY_j\right)
 =\KK_i\left(\bigoplus_j C_0(Y_j),\C\right)
 \cong\prod_j K_i^{\lf}(Y_j).
\end{equation}
This is an integral product. Restriction to a clopen union of components is
the corresponding coordinate projection, induced contravariantly by
extension by zero on $C_0$.

For an oriented circle $C$, choose its analytic fundamental class $[C]$ so
that its index pairing with a unitary of winding one is $1$. The circle index
formula then reads
\begin{equation}\label{eq:winding}
 \langle[f],[C]\rangle=\deg(f),\qquad f:C\to S^1.
\end{equation}
This fixes the otherwise harmless sign convention for the Dirac or Toeplitz
description of the class \cite[Example~4.2.4]{Khalkhali}.
An integral simplicial $1$-cocycle on a finite complex $Y$ determines a class
in $H^1(Y;\Z)$, represented by a map $f:Y\to S^1$, and hence a unitary
class $[f]\in K_1(C(Y))$. For an included oriented circle, naturality of the
index pairing identifies its pairing with the sum of the cocycle around
that circle \cite[Proposition~18.7.1 and Corollary~18.10.3]{Blackadar}.
One can see the representing map directly: map vertices to
$1$, traverse each edge with the assigned integer winding, and extend over
triangles because their winding sums vanish. Higher extensions exist since
$\pi_k(S^1)=0$ for $k\ge2$.
